\section{Discussion}
\label{sec:discussion_dynamic}

\subsection{What is explained that existing theories do not explain?}

The principal advance is not another instantaneous instability criterion. Existing scalar thermoplastic theories explain how hardening, rate sensitivity, thermal softening, inertia and diffusion enter a frozen dispersion relation \cite{Bai1982,WrightBatra1985,MolinariClifton1987,Wright2002}. Frozen crystal tangents add orientation and slip-system resolution. Neither construction, by itself, answers the history-selection question posed here: among all admitted crystal-resolved initial disturbances, which direction--wavenumber--state combination gains most over a specified interval while the base state changes?

The same-branch discrimination makes this distinction falsifiable. If accumulated frozen spectral abscissae were an adequate surrogate, $\Delta_{\rm FT}^{\rm full}$ in Eq.~\eqref{eq:ft_discrimination} would remain small when the base path, branch, 69-state space and norm are locked. Instead, the like-for-like discrepancies are $+4.784$, $+7.217$ and $+6.402$. The previously apparent negative onset-$x$ discrepancy arose from subtracting a full-state frozen rate from a projected constitutive gain; separating those two input--output questions removes that artefact without removing the non-equivalence. No branch-independent shift or multiplicative calibration of the frozen rate recovers the three propagated histories.

The propagator resolves the missing chronological information. It maps every admissible initial disturbance to the terminal observed state, and its leading singular triplet identifies the initial composition, cumulative gain and terminal composition that are optimal for the declared norm and window. The measured leading-mode condition numbers reach $2.96\times10^5$, adjacent eigenvectors rotate by as much as $51.1^\circ$, and adjacent generators do not commute. On the terminal branch, reversing the chronological factors lowers the log gain from 13.370 to 10.979, while the first-Magnus surrogate gives 10.115. These diagnostics identify non-orthogonal modal geometry, rotating instantaneous subspaces and operator ordering as active sources of the frozen/propagated difference. Optimizing the resulting map over projective direction and wavenumber produces a near-optimal \emph{set}, not merely a largest root. The theory therefore explains three observations unavailable from a single frozen spectrum: why two different directions are effectively tied near onset, why one branch later separates by orders of magnitude, and why the selected wavenumber drifts from a short onset scale to a longer terminal scale.

\subsection{A mechanism is a history-dependent coupling architecture}

Labeling the dominant terminal singular-vector block as ``plastic'' is descriptive but incomplete. The intervention results show why. At onset, dislocation--plastic coupling carries the largest absolute Shapley attribution; at the terminal horizon, inertia--plastic coupling does. The full attribution vectors have low cosine similarity and one cross-edge changes sign. Thus the terminal plastic state can be produced through different upstream coupling architectures at different horizons. This resolves an ambiguity common in localization discussions: a coordinate with the largest amplitude need not be the coupling whose removal changes amplification most.

The fixed-base interventions are intentionally more modest than causal identification in a physical specimen. Removing a Jacobian block changes the linearized model while holding its base path fixed; it answers a counterfactual question about the tangent architecture, not how a material would re-equilibrate after a constitutive mechanism were physically removed. Nevertheless, recomputation of all 64 coalitions is stronger evidence than reading components of one eigenvector or one singular vector, and it demonstrates temporal reorganization without conflating it with base-state changes.

\subsection{Relation to Bai theory and crystal plasticity}

The framework should be read as a crystal-resolved generalization of the \emph{question} posed by Bai, not yet as a completed asymptotic reduction to Bai's coefficients. The scalar theory and the HCP descriptor have both been independently audited. What remains open is a controlled chain of assumptions---single effective slip, isotropic/specified orientation reduction, frozen internal variables, compatible thermal variables and elimination order---under which the HCP determinant reduces analytically to the classical scalar polynomial. Until that chain is written and verified, ``recovers Bai exactly'' would be stronger than the evidence.

This boundary does not weaken the finite-time result. The core discrimination is internal to the HCP model: the frozen and propagated calculations use the same generator, complete state space and scaling; they differ only in whether chronological operator composition is retained. Therefore the conclusion that frozen accumulation is not equivalent does not depend on a cross-model Bai calibration. Bai remains the classical anchor; the new object is the crystal-resolved nonautonomous propagator.

\subsection{How far does the theory close cross-scale length and energy?}

The present theory closes one important part of the cross-scale problem. Crystal orientation, slip-system kinetics, mobile/dipole dislocation evolution, heat generation, inertia, conduction and gradient recovery all enter one finite-time operator. An interior $k^*$ for a constitutive-response selector is a finite-band selection result, and the first-law ledger partitions crystal-plastic power into heat and a passive nonthermal remainder without identifying the latter as dislocation free energy. The unrestricted full-state statement is weaker: its onset optimum is a high-$k$ boundary wave that continues to grow when the compact search interval is extended. The selector-specific re-optimization therefore identifies exactly which input--output questions admit a finite localization wavelength and which do not.

It does \emph{not} yet close the full chain from grain/void scale to process-zone width, propagation resistance and structural band thickness. Four gaps remain. First, $2\pi/k^*$ is the period of an optimally amplified linear disturbance, not the width of a mature localization zone. Second, the separate same-window width audit finds that accepted single-mode widths scale with the imposed wavelength and that multimode widths are under-resolved; no seed-independent material width emerges. Third, neither an identified internal-variable Helmholtz potential nor a propagation-resistance or fracture-like energy per unit area has been derived from the volumetric power ledger. Fourth, although a four-mode Galerkin replay reproduces the optimized linear inputs, retaining every Fourier mode resolvable on the same 16-cell grid makes the nonlinear continuation non-finite even after substep refinement. Thus the current crystal-plasticity--perturbation theory closes finite-time finite-band \emph{linear selection}; it does not close nonlinear harmonic transfer, finite width, finite dissipation or propagation resistance across scales.

The negative width result is consequential. Earlier cross-horizon FWHM/wavelength ratios can organize numerical cases, but they do not identify a constitutive transfer coefficient because their windows differ. Even under a matched window, grid convergence alone is insufficient when the converged width tracks the imposed seed. A future material-length claim must pass the complete denial gate: broadband/multimode initiation, phase and spectrum insensitivity, at least 12 cells across the peak, refinement stability, and an energy balance that yields a finite areal dissipation or propagation resistance.

\subsection{Mathematical and evidential limits}

The critical time in Eq.~\eqref{eq:ft_result_tc} is a threshold crossing in a declared dimensionless norm. Equivalent finite-dimensional norms preserve finite-time boundedness, but they do not preserve a numerical $\mathcal G=e$ crossing. The complete 11-by-7 fixed-candidate screen is therefore only a sensitivity map over three registered branches. The unrestricted 70-problem re-optimization supplies the stronger result: every constitutive/thermal-response selector retains an interior optimum under the five audited norms, while full/full and mechanical/mechanical can select a high-$k$ boundary wave. This is robustness conditional on a declared physical question, not coordinate or selector invariance. None of the 60 non-rank-one optimized maps triggers the formal singular-value degeneracy gate; the ten scalar-temperature maps have no $\sigma_2$. This pointwise separation does not remove the set-valued near-onset competition between distinct direction--wavenumber basins.

The strengthened search combines projective-sphere--$\log k$ reconnaissance, exact symmetry and prior-branch anchors, basin recurrence, boundary classification and 129-state reevaluation. Because blind Sobol sampling can miss a narrow basin, it remains an anchor-assisted basin-retention and convergence audit, not an independent global-search certificate or a proof over every basin in $\mathbb{RP}^2\times[k_{\min},k_{\max}]$. The strictly same-selector piecewise-frozen control converges under interval-preserving refinement; this verifies the propagation implementation but does not turn the scalar frozen-abscissa integral into an equivalent vector propagator.

Positive sampled acoustic, conduction and gradient symbols remove obvious pathologies but do not prove a uniform high-$k$ bound for a non-normal propagator. Proposition~5 states the missing symmetrizer condition explicitly. The current high-$k$ calculations are consistency evidence on registered paths, not construction of that symmetrizer. Constitutive switching surfaces are excluded because ordinary branch-locked Jacobians are used; a semismooth or measure-differential extension would be required to cross them rigorously.

The failed full-Fourier continuation is consistent with this limitation and
prevents using the low-mode replay as a substitute for a high-$k$ theorem.
It does not by itself distinguish physical short-wave instability from loss
of continuum validity, aliasing or insufficient nonlinear regularization.
Closing that distinction requires a convergent de-aliased spatial ladder and
a constitutively justified high-gradient or nonlocal energy, rather than an
ad hoc spectral cutoff.

Finally, the titanium-like card exercises the complete operator but is not identified against a matched batch, loading history and full-field localization measurement. The direct comparison with synchronized Grade-II CP-Ti data makes the limitation quantitative: at a matched nominal rate the model peaks far too early, produces almost no temperature rise at that peak and cannot reproduce the post-localization stress collapse. Appendix~\ref{app:parameter_provenance} makes the provenance distinction parameter by parameter: heat-capacity and selected thermal data have direct reference provenance, the crystal/slip block is literature constrained but largely inherited from a watermarked verification seed, and the gradient lengths are registered model controls rather than measured band widths. Titanium experiments demonstrate that shear response, GND accumulation, alloy state and band self-organization provide material-specific constraints that a verification card cannot replace \cite{GuoEtAl2019PRL,ZhuEtAl2017,RanEtAl2021,XueMeyersNesterenko2002}. The predicted $t_c$, gain and $k^*$ should therefore be interpreted as theory-discriminating computational observables. Specimen-level prediction requires independent calibration and validation after, not before, the operator and scale gates are closed.
